\section{Introduction}\label{sec:introduction}

Let \(Y\) be a smooth projective complex variety.  Iritani's Gamma integral
structure associates a flat section \(s_Y(E)\) of the quantum connection to a
topological \(K\)-class \(E\) and identifies the flat pairing with the Euler
pairing \cite{IritaniGamma}.  If \(p\in Y\) is a point, then
\[
 [s_Y(E),s_Y(\OO_p))=\chi(E,\OO_p)=\rk(E).              \tag{1.1}
\]
The second entry in (1.1) is therefore a distinguished flat covector.  We
call it the \emph{Gamma point row}.

We study whether the zero or nonzero restriction of this row to a formal
monodromy packet survives birational modification.  A direct composition of
one-wall quantum comparisons does not answer the question: adjacent walls
use incompatible Novikov completions, and changing a sectorial lift can add
a rank-visible Stokes correction.  The main construction instead puts the
two birational endpoints in one equivariant cobordism and compares their
localized gauged theories before either exceptional completion is taken.

The resulting conditional invariant gives a precise route to stable
irrationality for every smooth cubic threefold.

\begin{theorem}[Conditional cubic stabilization criterion]
\label{thm:intro-cubic-conditional}
Let \(X\subset\mathbf P^4_{\C}\) be a smooth cubic threefold.  Assume that
every smooth projective birational map admits a W{\l}odarczyk completion
satisfying tailwise derived identification and marked threshold compatibility,
in the sense of Hypotheses~\ref{hyp:tailwise-derived} and
\ref{hyp:marked-threshold}.  Then
\(X\times\mathbf P^m\) is irrational for every \(m\geq0\).
\end{theorem}

The \(X\times\mathbf P^1\) case has a separate
dimension-four proof by weak factorization \cite{RuddEpilogue}.  We keep that
argument separate: the present proof introduces a global gauged-localization
mechanism which is independent of the stabilization dimension.  The case
\(m=2\) is the first one not covered by the dimension-four argument.

The conditional conclusion is independent of the classical decomposition of
the diagonal and integral Hodge shadows of stable rationality.  The companion paper
exhibits universally \(\mathrm{CH}_0\)-trivial cubic threefolds, for which that
cycle-theoretic obstruction vanishes, while their first stabilizations remain
irrational \cite{RuddEpilogue}.  The present invariant records irregular
quantum-monodromy data instead; it makes no assertion about the Hodge
conjecture.

The unconditional endpoint contrast is
\[
 b(X\times\mathbf P^m,\cP_6)=1,
 \qquad b(\mathbf P^{m+3},\cP_6)=0.                          \tag{1.2}
\]
where an empty packet has value zero.  Here \(b\) records whether the point
row is nonzero on the indicated packet.  Section~\ref{sec:cubic-endpoint}
proves (1.2).  Cai's cubic connection supplies the primitive-sixth block; we
reconstruct its indicial polynomial from his matrices and compute the point
coefficient independently from the cubic hypergeometric period.  Thus the
cubic-specific source dependency is confined to one endpoint lemma.

The birational statement behind the application is more general.

\begin{theorem}[Conditional point-row primary invariance]
\label{thm:intro-birational-conditional}
For smooth projective birational complex varieties whose global cobordism
satisfies Hypotheses~\ref{hyp:tailwise-derived} and
\ref{hyp:marked-threshold}, the Gamma point row is nonzero on a
formal-monodromy primary packet at one endpoint if and only if it is nonzero
on the corresponding packet at the other endpoint.
\end{theorem}

The conditional proof uses a smooth projective equivariant completion of one
W{\l}odarczyk cobordism.  A point in the trivialized common birational open
sweeps out an orbit
cylinder.  Its equivariant class restricts to the two endpoint point classes
and vanishes on all intermediate fixed strata.  Rotation localization
therefore removes every wall contribution.  For each neutral affine gauge
direction, Woodward's factorization packages the complete bubble tree into
fixed stable-map factors.  Hypothesis~\ref{hyp:tailwise-derived} identifies
the fixed stacks, perfect obstruction theories, virtual classes, and
evaluations along each affine tail; the remaining moving factors then give a
finite-jet Gamma recurrence, so each tail is holonomic and tempered.
Hypothesis~\ref{hyp:marked-threshold} asks that the finite cyclic Rees
\(z\)-modules be identified across the intervening thresholds, with formal
monodromy and the marked point row preserved.  Nonneutral directions are
Laurent-finite, and Rees homogenization makes both endpoint
maps intertwine formal monodromy with the same half-Tate shift.  The dual
cyclic module generated by the continued row then determines exactly which
primary packets it detects.  This is weaker than continuing the complete
nonlinear quantum source.  It remains an open marked connection problem, not
a consequence of the available linear GLSM contour theorem.

The paper also records two local results which motivated the global
construction.  The first concerns a smooth projective simple VGIT wall
\(X_-\dashrightarrow X_+\) in the sense of Gu--Yu--Yu, oriented so that
\(r_+<r_-\) \cite{GuYuYu}.  Their comparison is a formal, pairing-preserving
decomposition
\[
 \QDM(X_-)
 \cong \QDM(X_+)\oplus\bigoplus_j\QDM(S)_j               \tag{1.3}
\]
over an exceptional Laurent completion, where \(S\) is the wall quotient.
Choose corresponding points \(p_\pm\) in the common open set.

\begin{theorem}[Common-open point column]
\label{thm:intro-simple-wall}
At the extremal specialization used in the Gu--Yu--Yu comparison, the
ambient coordinate of the point class is exact when \(r_+<r_-\):
\[
 \pr_{X_+}\Phi([p_-])=[p_+].                              \tag{1.4}
\]
No assertion is made that every wall coordinate of \(\Phi([p_-])\) vanishes.
\end{theorem}

The proof uses the distinguished master-space lift of the common point.  Its
wall restriction vanishes.  Applying the negative-chamber Fourier
isomorphism to the lift and to its wall-frequency derivative shows that the
derivative is zero in the specialized completed source.  Fourier covariance
then makes the positive point image horizontal.  A regular-singular
recursion kills its apparent positive wall-parameter tail.

Our second result is analytic but crepant.  Let \(Y\dashrightarrow Y'\) be a
projective ordinary flop covered by Lee--Lin--Qu--Wang
\cite{LeeLinQuWang}.  Fix one of their analytic-continuation domains in the
extremal variable.

\begin{theorem}[Ordinary-flop point row]
\label{thm:intro-flop}
The analytically continued point section of \(Y\) equals the intrinsic point
section of \(Y'\) on the fixed continuation domain.  Consequently the point
row has the same zero or nonzero restriction to every sectorially realized
formal-monodromy packet whose chosen realization is transported by the graph
gauge on that branch.
\end{theorem}

Pure extremal curves lie in the exceptional locus and cannot meet a point in
the common open set.  This computes the point column exactly at transverse
Novikov degree zero.  A divisor pulled back from the common contraction then
kills the transverse tail coefficient by coefficient.  This is a theorem
about one distinguished column; it does not promote the ancestor comparison
of \cite{LeeLinQuWang} to full descendant invariance.

The discrepant local case contains a further frame comparison.  A formal summand,
a sectorial lift of that summand, and the intrinsic large-radius Gamma
section are different objects.  Gu--Yu--Yu prove (1.3) in a Laurent/formal
coefficient ring, while Shen--Shoemaker identify Gamma asymptotic classes at
an extremal specialization \cite{ShenShoemaker}.  Passing from these results
to an intrinsic Gamma point row requires a common pairing-preserving
sectorial realization.  We state the corresponding one-wall consequence
only under that explicit hypothesis.

The negative results show why this distinction is necessary.  A rank-two
incomplete-Gamma connection has a nonzero Stokes shear from a
\(\zetaSix\)-formal line into a rank-visible ambient line.  The example admits
a flat integral hyperbolic doubling and remains after tensoring with
\(\Q[N]/(N^3)\).  Localized partial Fourier--Laplace transform kills
the constant extension which records the shear.  In two Fourier variables,
the delta module at the common origin is invisible after either localization
but transforms back to a constant full-support module.  Boundary support in
Fourier coordinates therefore points in the opposite direction from the
desired rank conclusion.

At the local level, the surviving two-wall statement is one-sided.  Complete ray-ordered
corrections must have rank-zero \emph{targets}; their sources and scalar
coefficients are irrelevant.  Under that hypothesis, a finite factorization
preserves the point row by elementary linear algebra.  A signed punctual
coefficient of an oriented boundary complex is a plausible shadow of these
factorwise conditions, but it is not equivalent to them in the local
formalism.  A single alternating coefficient can vanish by cancellation.
Under Hypotheses~\ref{hyp:tailwise-derived} and
\ref{hyp:marked-threshold}, the global construction in
Section~\ref{sec:global-transport} bypasses this missing marked comparison.

Sections~\ref{sec:point-row}--\ref{sec:ordinary-flop} establish the point-row
formalism and the local identities.  Sections~\ref{sec:incomplete-gamma} and
\ref{sec:fourier-boundary} give the countermodels, and
Section~\ref{sec:two-wall} isolates the local composition condition they
force.  Section~\ref{sec:global-transport} proves the Gamma-ratio reduction
and conditional Theorem~\ref{thm:intro-birational-conditional} by the
common-master method.  Section~\ref{sec:cubic-endpoint} proves (1.2) and
Theorem~\ref{thm:intro-cubic-conditional}.  Section~\ref{sec:scope} records the source,
analytic, and verification boundaries.
